\section{The Church at Worship: Water, Bread, Prayer, and Custom}

\subsection{A baptismal body}

\emph{On Baptism} is the earliest surviving Latin treatise devoted to baptism.
It answers Christians unsettled by the apparent simplicity of water and by a
woman teacher whom Tertullian calls a viper of the Cainite heresy. The insult
belongs to his polemic; the existence of doctrinal argument around baptism is
the historical setting. Chapters 3--5 move through creation, the Spirit over
the waters, pagan washings, and biblical types. The material element is not too
low for God. The whole treatise anticipates his larger defense of flesh.

Chapter 4 describes prayer invoking God over the water and the Spirit's
sanctifying action. Chapters 6--8 move from profession under the threefold name
through emergence from the water, anointing, imposition of the hand, and
invocation or reception of the Holy Spirit. Chapter 13 relates the command of
Matthew 28:19 to the Trinitarian name. Tertullian's arrangement is a precious
witness to African Christian initiation, but his theological exposition is not
a verbatim service book. Words, ministers, and order must not be standardized
beyond what he actually reports.

His practical judgments are equally revealing. Chapter 17 assigns ordinary
authority to baptize to the bishop, then presbyters and deacons with his
permission, while allowing lay baptism in necessity. Tertullian immediately
restricts that concession: emulation of episcopal office breeds schism, and
women may not baptize. He rejects an appeal to the \emph{Acts of Paul and
Thecla} by reporting that an Asian presbyter confessed to composing the work
out of love for Paul and was removed. The anecdote is early reception of that
text; it is not surviving trial minutes and does not settle every part of the
work's composition.

Chapter 18 counsels delay where candidates lack maturity, explicitly mentioning
little children and unmarried people vulnerable to temptation. The passage is
evidence that the baptism of children was known and debated; it is not proof
that all African churches delayed it or that Tertullian denied baptismal grace.
His fear is grave post-baptismal sin. Chapter 19 prefers Passover and the period
through Pentecost as especially fitting times but states that every day belongs
to the Lord and baptism may be celebrated at any time. Chapter 20 calls for
prayer, fasting, vigils, confession, and humility before and after the water.

Chapter 16 adds ``baptism of blood'': martyrdom can supply the washing not
received or restore what was lost. The language binds sacrament, confession,
and body; it also helps explain why compromised post-baptismal life appeared so
terrible to him. What persecution could seal at the cost of blood, ordinary
leniency must not cheapen.

\subsection{Inherited actions and unwritten reasons}

\emph{On the Crown} 3 gives a compact list of Christian practices: renunciation
of the devil before baptism, triple immersion, responses somewhat fuller than
the Gospel wording, tasting milk and honey, abstaining from the daily bath for
a week, receiving the Eucharist before dawn and only from the hands of
presiders, annual offerings for the dead, refusal to kneel on Sunday and during
Pentecost, and tracing the sign of the cross repeatedly in daily life. Chapter
4 says no written law from Scripture can be produced for every such act; custom
confirmed by tradition and supported by faith governs them.

The passage became a classic witness for tradition. It must be read at the
scale of its argument. Tertullian is defending an unwritten refusal to wear a
military crown by analogy with recognized ecclesial customs. He establishes
that these practices were persuasive examples in his circle. He does not
provide the date of their origin, prove that every church practiced every item,
or grant later customs automatic immunity from judgment. Nor does his appeal to
tradition cancel his frequent demand that customs be tested by truth.

The noun \emph{sacramentum} in Tertullian can mean oath, sacred bond, divine
plan, or a Christian rite according to context. Its use helped furnish later
Latin sacramental vocabulary, but it is anachronistic to import a complete
medieval sacramental system into every occurrence. His evidence is especially
valuable when work, locus, and action are named rather than when one technical
word is made to carry centuries of development.

\subsection{Prayer through the day}

\emph{On Prayer} 1 calls the Lord's Prayer a summary of the whole Gospel. The
treatise moves clause by clause and then into bodily and communal practice.
Chapter 6 reads the petition for daily bread both ordinarily and with reference
to Christ and the Eucharist. Chapter 23 defends standing rather than kneeling on
the Lord's Day and during the paschal season; chapter 25 recognizes customary
hours of prayer, especially the third, sixth, and ninth, while making prayer at
the beginning and end of day fundamental. Chapter 27 mentions a psalm after
common prayer. Chapters 28--29 present prayer from a chaste body and innocent
soul as the spiritual sacrifice God desires.

These instructions make early Christian prayer embodied: washed hands do not
replace moral innocence; an ostentatious voice does not compel God; dress,
kiss, posture, time, psalm, and household relation all matter. They also expose
contested practice. The treatise is an argument within worship, not a camera
placed unnoticed in a universally uniform liturgy.

\subsection{Bread, body, and a prophet's report}

Against Marcion, Tertullian insists that Christ at the supper made the bread his
body, describing it as a figure of his body (\emph{Against Marcion} 4.40). The
sentence has been enlisted in later contests between ``symbol'' and ``real
presence.'' Tertullian's target is docetism: a figure of the body presupposes a
true body capable of being represented. His Eucharistic sign and his realism
work together. The passage is an anti-Marcionite argument, not a complete
account of later Catholic Eucharistic doctrine.

\emph{On the Soul} 9 preserves a rarer scene. A Christian woman endowed with
revelatory gifts receives visions during the Sunday rites amid Scripture,
psalms, preaching, and prayer; after the people are dismissed, what she saw is
reported and examined. Tertullian uses her vision as evidence for the soul's
corporeality. A visionary report cannot prove the ontology he derives from it.
Historically, however, the scene shows prophecy neither wholly private nor
simply sovereign: it occurs within assembly and is tested afterward. The woman
whose speech supplies theological evidence will complicate Tertullian's severe
limits on women's teaching and ministry.

Across these works the Church is not an abstraction. It washes, anoints,
imposes hands, tastes, eats, stands, kneels, fasts, sings, signs the body, tests
speech, gives alms, and remembers the dead. That density is one reason
Tertullian survived even where his later conclusions did not. He makes ancient
Christian life visible precisely because every bodily practice can become an
argument.
